\begin{lemma}
Certified ordered tree-sequences transport in both directions between $G$
and $G^{\circ}$ without changing an edge identity.  More precisely, for
every $e_0$, palette size $k$, coloring $\varphi$, and ordered tree-sequence
$T$ of $G$, there is an ordered tree-sequence $T^{\circ}$ of $G^{\circ}$
with the same length, the same edge at every ordered edge position, and at
every ordered vertex position the incident-core copy of the corresponding
vertex of $T$.  Their vertex sets have equal cardinality, and
$T^{\circ}$ is a Tashkinov tree exactly when $T$ is.  Conversely, coercing
the vertices of any ordered core tree-sequence to $V(G)$ gives such a
sequence of $G$ with the same indexed edges and vertices, equal vertex-set
cardinality, and equivalent Tashkinov-tree certification.
\end{lemma}
\begin{proof}
We first transport a tree-sequence $T$ of $G$ to the core.  Every indexed
vertex of $T$ is incident with an indexed edge.  Indeed, the vertex in
position $r+1$ is an endpoint of the edge in position $r$ by the attachment
clause of \noderef{TreeSequence}.  For the vertex in position $0$, use the
edge in position $0$, which exists because the length is positive: its
attachment index lies in the one-element initial segment and therefore is
$0$.  Thus each indexed vertex has a witness placing it in the subtype
\noderef{IncidentVertex}.

Keep the length and every edge entry unchanged and replace each vertex entry
by this incident-core copy.  Injectivity of the vertex list follows from
injectivity of the original list after forgetting subtype witnesses;
injectivity of the edge list and the first-edge equation are literally
unchanged.  For an attachment equation, apply
\noderef{IncidentCoreEndpoints}: the coerced core endpoint pair is the
original endpoint pair.  Both displayed endpoints are the incident-core
copies of the original displayed endpoints, and subtype equality is
determined by the underlying vertex.  Hence the original attachment equation
lifts to the required core attachment equation.  This constructs the
ordered core sequence with exactly the stated indexed vertices and edges.
The indexed vertex bijection gives equality of the two vertex-set
cardinalities in \noderef{TreeSequenceVertexSet}.

To compare the Tashkinov certificates, fix a positive edge position $j$.
The edge identity and its assigned color are unchanged.  A candidate earlier
vertex has the same index on both sides, and
\noderef{IncidentCoreMissingColors} says that precisely the same colors are
missing at that vertex and at its underlying original vertex.  Therefore the
existential earlier-index and missing-color witness required by
\noderef{TashkinovTree} transports in each direction.  Quantifying over $j$
proves the asserted equivalence of Tashkinov-tree certification.

Conversely, start with a core tree-sequence and coerce every indexed vertex
to $V(G)$, retaining the same length and edge list.  Vertex injectivity is
preserved because equality of underlying vertices implies equality in the
incident-vertex subtype.  The first-edge equation is unchanged, and applying
\noderef{IncidentCoreEndpoints} to each core attachment equation gives the
original attachment equation.  The indexed vertex map is again a bijection
between the two vertex lists, hence preserves vertex-set cardinality.  The
same missing-color argument using
\noderef{IncidentCoreMissingColors} proves the Tashkinov equivalence.  This
establishes all stated data in both directions, with every parallel edge
identity kept at its original ordered position.
\end{proof}
